\subsection{INTEGRATION OF SECTIONS OF A VECTOR BUNDLE}

In this number and the next, we denote by X a real manifold of class \(C^{r}\) \((1 \leq r \leq \infty)\) , locally of finite dimension.

Let Y be a manifold of class \(C^{r}\) . If \(r < \infty\) , consider the map \(f \mapsto j^{r}(f)\) from \(\mathcal{C}^{r}(\mathrm{X};\mathrm{Y})\) to \(\mathcal{C}(\mathrm{X};\mathrm{J}^{r}(\mathrm{X},\mathrm{Y}))\) (Differentiable and Analytic Manifolds, Results, 12.3.7). The inverse image under this map of the topology of compact convergence on \(\mathcal{C}(\mathrm{X};\mathrm{J}^{r}(\mathrm{X},\mathrm{Y}))\) is called the topology of compact \(C^{r}\) -convergence on \(\mathcal{C}^{r}(\mathrm{X};\mathrm{Y})\) ; it is the upper bound of the topologies of uniform \(C^{r}\) -convergence on K (Differentiable and Analytic Manifolds, Results, 12.3.10), where K runs through the set of compact subsets of X.

When \(r = \infty\) , we call the topology of compact \(C^{\infty}\) -convergence on \(\mathcal{C}^{\infty}(X;Y)\) the upper bound of the topologies of compact \(C^{k}\) -convergence, in other words the coarsest topology for which the canonical injections \(\mathcal{C}^{\infty}(X;Y) \to \mathcal{C}^{k}(X;Y)\) are continuous for \(0 \leq k < \infty\) .

Let \(\mathrm{E}\) be a real vector bundle with base \(\mathrm{X}\) , of class \(\mathrm{C}^r\) , and let \(\mathcal{S}^r(\mathrm{X};\mathrm{E})\) be the vector space of sections of \(\mathrm{E}\) of class \(\mathrm{C}^r\) . In this number we give \(\mathcal{S}^r(\mathrm{X};\mathrm{E})\) the topology induced by the topology of compact \(\mathrm{C}^r\) -convergence on \(\mathcal{C}^r(\mathrm{X};\mathrm{E})\) , also called the topology of compact \(\mathrm{C}^r\) -convergence; it makes \(\mathcal{S}^r (\mathrm{X};\mathrm{E})\) into a complete separated locally convex topological vector space (cf.~Differentiable and Analytic Manifolds, Results, 15.3.1 and Spectral Theories, in preparation).

Now let H be a Lie group, \(m: H \times X \to X\) a law of left operation of class \(C^{r}\) ; put \(hx = m(h, x)\) for \(h \in H\) , \(x \in X\) . Let E be a vector H-bundle with base X, of class \(C^{r}\) (Chap. III, §1, no. 8, Def. 4). For \(s \in \mathcal{S}^{r}(X; E)\) and \(h \in H\) , denote by \(^{h}s\) the section \(x \mapsto h.s(h^{-1}x)\) of E; the map \((h, s) \mapsto ^{h}s\) is a law of operation of H on the space \(\mathcal{S}^{r}(X; E)\) .

Lemma 4. The law of operation \(\mathrm{H} \times \mathcal{S}^r (\mathrm{X};\mathrm{E}) \to \mathcal{S}^r (\mathrm{X};\mathrm{E})\) is continuous.

In view of the definition of the topology of \(\mathcal{S}^{r}(\mathrm{X};\mathrm{E})\) and General Topology, Chap. X, §3, no. 4, Th. 3, it suffices to prove that for any integers \(k \leq r\) , the map \(f : H \times X \times \mathcal{S}^{k}(\mathrm{X};\mathrm{E}) \to \mathrm{J}^{k}(\mathrm{X};\mathrm{E})\) such that \(f(h,x,s) = j_{x}^{k}(^{h}s)\) is continuous. For \(h \in H\) , denote by \(\tau_{h}\) (resp. \(\theta_{h}\) ) the automorphism \(x \mapsto hx\) of X (resp. of E). Define maps

\[
f _ {1}: \mathrm{H} \times \mathrm{X} \to \mathrm{J} ^ {k} (\mathrm{X}, \mathrm{X})
\]

\[
f _ {2}: \mathrm{H} \times \mathrm{E} \rightarrow \mathrm{J} ^ {k} (\mathrm{E}, \mathrm{E})
\]

\[
g: \mathrm{H} \times \mathrm{X} \times \mathcal {S} ^ {k} (\mathrm{X}; \mathrm{E}) \to \mathrm{J} ^ {k} (\mathrm{X}, \mathrm{E})
\]

by \(f_{1}(h,x) = j_{x}^{k}(\tau_{h}),f_{2}(h,v) = j_{v}^{k}(\theta_{h}),g(h,x,s) = j_{hx}^{k}(s)\) . We have

\[
f (h, x, s) = f _ {2} (h, s (h ^ {- 1} x)) \circ g (h ^ {- 1}, x, s) \circ f _ {1} (h ^ {- 1}, x),
\]

and consequently it suffices, by Differentiable and Analytic Manifolds, Results, 12.3.6, to prove that \(f_{1}, f_{2}\) and \(g\) are continuous.

Now \(g\) is the composite map

\[
\begin{array}{c c c} \mathrm{H} \times \mathrm{X} \times \mathcal {S} ^ {k} (\mathrm{X}; \mathrm{E}) & \xrightarrow {(m , \mathrm{Id})} & \mathrm{X} \times \mathcal {S} ^ {k} (\mathrm{X}; \mathrm{E}) \\ & \xrightarrow {(\mathrm{Id} , j ^ {k})} & \mathrm{X} \times \mathcal {C} (\mathrm{X}; \mathrm{J} ^ {k} (\mathrm{X}, \mathrm{E})) \quad \xrightarrow {\varepsilon} \quad \mathrm{J} ^ {k} (\mathrm{X}, \mathrm{E}) \end{array}
\]

with \(\varepsilon(x, u) = u(x)\) ; the map \(\varepsilon\) being continuous (General Topology, Chap. X, §3, no. 4, Cor. 1 of Th. 3), \(g\) is continuous.

Let \((h_{0}, x_{0}) \in \mathrm{H} \times \mathrm{X}\) ; we shall prove that \(f_{1}\) is continuous at \((h_{0}, x_{0})\) . There exist charts \((\mathrm{U}, \psi, \mathrm{F})\) and \((\mathrm{V}, \chi, \mathrm{F}')\) of \(\mathrm{X}\) and an open subset \(\Omega\) of \(\mathrm{H}\) such that \(x_{0} \in \mathrm{U}, h_{0} \in \Omega\) and \(m(\Omega \times \mathrm{U}) \subset \mathrm{V}\) . By using the expression for \(\mathrm{J}^{k}(\mathrm{X}, \mathrm{X})\) in these charts, we are reduced to proving, for \(1 \leq l \leq k\) , the continuity at \((h_{0}, x_{0})\) of the map \((h, x) \mapsto \Delta_{x}^{l}(\tau_{h})\) from \(\Omega \times \mathrm{U}\) to \(\mathrm{P}_{l}(\mathrm{F}; \mathrm{F}')\) , with \(\Delta_{x}^{l}(\tau_{h})(v) = \frac{1}{l!}\mathrm{D}^{l}\tau_{h}(x).v\) for \(v \in \mathrm{F}\) (Differentiable and Analytic Manifolds, Results, 12.2). But \(\mathrm{D}^{l}\tau_{h}(x)\) is simply the \(l\) th partial derivative of \(m(h, x)\) with respect to \(x\) , which is continuous by hypothesis; consequently, \(f_{1}\) is continuous. The proof that \(f_{2}\) is continuous is similar, hence the lemma.

PROPOSITION 5. Assume that the group H is compact and denote by dh the Haar measure on H of total mass 1. Let s be a section of E of class \(C^{r}\) .

Denote by \(s^{\sharp}\) the vector integral \(\int_{\mathrm{H}}{}^{h}s\,dh\) . Then \(s^{\sharp}\) is a section of \(\mathrm{E}\) of class \(\mathrm{C}^r\) , invariant under \(\mathrm{H}\) ; for \(x \in \mathrm{X}\) , we have \(s^{\sharp}(x) = \int_{\mathrm{H}}hs(h^{-1}x)\,dh \in \mathrm{E}_x\) . The endomorphism \(s \mapsto s^{\sharp}\) of \(\mathcal{S}^r(\mathrm{X};\mathrm{E})\) is a projection onto the subspace of H-invariant sections.

Consider the map \(h \mapsto {}^{h}s\) from H to \(\mathcal{S}^{r}(\mathrm{X};\mathrm{E})\) ; it is continuous by Lemma 4. Since the space \(\mathcal{S}^{r}(\mathrm{X};\mathrm{E})\) is separated and complete, the integral \(s^{\sharp} = \int_{H}{}^{h}s\,dh\) belongs to \(\mathcal{S}^{r}(\mathrm{X};\mathrm{E})\) (Integration, Chap. III, §3, no. 3, Cor. 2). The linear map \(s \mapsto s(x)\) from \(\mathcal{S}^{r}(\mathrm{X};\mathrm{E})\) to \(E_{x}\) being continuous, we have \(s^{\sharp}(x) = \int_{H}{}^{h}s(x)\,dh\) for all \(x \in X\) . It is clear that \(s^{\sharp}\) is invariant under H; if s is an H-invariant section, we have \(s^{\sharp} = s\) , hence the last assertion.

COROLLARY 1. Let \(\mathrm{F}\) be a Banach space, \(\rho : \mathrm{H} \to \mathbf{GL}(\mathrm{F})\) an analytic linear representation, \(f \in \mathcal{C}^r(\mathrm{X};\mathrm{F})\) . For \(x \in \mathrm{X}\) , put

\[
f ^ {\sharp} (x) = \int_ {\mathrm{H}} \rho (h). f (h ^ {- 1} x) d h.
\]

Then \(f^{\sharp}\) is a morphism of class \(\mathbf{C}^r\) from \(\mathbf{X}\) to \(\mathbf{F}\) , compatible with the operations of \(\mathbf{H}\) ; for \(x \in \mathbf{X}\) , we have (with \(\tau_h\) denoting the automorphism \(x \mapsto hx\) of \(\mathbf{X}\) )

\[
d _ {x} f ^ {\sharp} = \int_ {\mathrm{H}} (\rho (h) \circ d _ {h ^ {- 1} x} f \circ \mathrm{T} _ {x} (\tau_ {h ^ {- 1}})) d h \in \mathscr {L} (\mathrm{T} _ {x} (\mathrm{X}); \mathrm{F}).\tag{14}
\]

The first assertion follows from the proposition applied to the bundle \(X \times F\) , equipped with the law of operation \((h; (x, f)) \mapsto (hx, \rho(h).f)\) . The second follows from Integration, Chap. III, §3, no. 2, Prop. 2, by applying to the vector integral \(f^{\sharp}\) the homomorphism \(d_{x} : \mathcal{C}^{r}(\mathrm{X};\mathrm{F}) \to \mathcal{L}(\mathrm{T}_{x}(\mathrm{X});\mathrm{F})\) which is continuous by the definition of the topology of compact \(C^{r}\) -convergence.

COROLLARY 2. Let F be a Banach space, \(f \in \mathcal{C}^{r}(X;F)\) ; put

\[
f ^ {\sharp} (x) = \int_ {\mathrm{H}} f (h x) d h
\]

for \(x \in \mathrm{X}\) . The function \(f^{\sharp}\) is of class \(\mathbf{C}^r\) , and \(f^{\sharp}(hx) = f^{\sharp}(x)\) for \(x \in \mathrm{X}\) , \(h \in \mathrm{H}\) .

COROLLARY 3. Let F be a Banach space, p an integer \(\geq 0\) , \(^{k}\Omega^{p}(\mathrm{X};\mathrm{F})\) the space of differential forms of degree p on X, with values in F, and of class \(C^{k}\) ( \(2 \leq k + 1 \leq r\) ). For \(\omega \in {}^{k}\Omega^{p}(\mathrm{X};\mathrm{F})\) , put \(\omega^{\sharp} = \int_{H} \tau_{h}^{*}\omega dh\) . Then the map \(\omega \mapsto \omega^{\sharp}\) is a projection on \(^{k}\Omega^{p}(\mathrm{X};\mathrm{F})\) whose image is the subspace of H-invariant forms. We have \(d(\omega^{\sharp}) = (d\omega)^{\sharp}\) for all \(\omega \in {}^{k}\Omega^{p}(\mathrm{X};\mathrm{F})\) .

The first assertion follows from the proposition applied to the vector H-bundle \(\operatorname{Alt}^{p}(\mathrm{T}(\mathrm{X});\mathrm{F})\) (Chap. III, §1, no. 8, Examples). To prove the second assertion, it suffices, in view of Integration, Chap. III, §3, no. 2, Prop. 2, to prove that the map \(d:{}^{k}\varOmega^{p}(\mathrm{X};\mathrm{F})\to{}^{k-1}\varOmega^{p+1}(\mathrm{X};\mathrm{F})\) is continuous when the first (resp. second) space is given the topology of compact \(C^{k}\) -convergence (resp. \(C^{k-1}\) -convergence). But this follows immediately from the definition of these topologies by means of semi-norms (Spectral Theories, in preparation) and the fact that d is a differential operator of order \(\leq 1\) (Differentiable and Analytic Manifolds, Results, 14.4.2).

